\section{Rank grading: cancellation and determinant families}
\label{sec:grading}

The strong-coupling coefficients of \S\ref{sec:second} and
\S\ref{sec:fourth} are exact rational functions of $N$. That makes a
question available which is invisible at any single rank: how does each
order of the series scale with the rank, and why?

The answer has two halves. At even orders the scaling is fixed by an
exact cancellation among four representation channels. At odd orders
the coefficients vanish outright for even $N$, and for odd $N$ they
first appear at order $N-2$ with a closed-form vertex. Neither statement
can be seen from numerics at $N = 3$.

\subsection{The one-face grading law}

Let $S_m(N)$ be the $C$-parity splitting of the one-plaquette
strong-coupling diagonal at order $m$,
\begin{equation}
  S_m(N) \;=\; (\text{odd})_m - (\text{even})_m \;=\; -2\,K_m[0][1].
\end{equation}

\begin{theorem}[Single-face planar grading through tenth order]
\label{thm:grading}
At even orders $m = 2, 4, 6, 8, 10$ the exact rational function of $N$
satisfies
\begin{equation}
  S_m(N) \;=\; c_m\, N^{-(2m-1)}\big(1 + O(N^{-2})\big),
  \label{eq:grading}
\end{equation}
with
\begin{equation}
  c_2 = -4, \quad c_4 = -32, \quad c_6 = -\tfrac{1748}{3},
  \quad c_8 = -\tfrac{123332}{9}, \quad c_{10} = -\tfrac{49593808}{135}.
\end{equation}
Writing $u = N^2\tau/2$ as in \eqref{eq:tau}, the \emph{leading} part of
the order-$m$ contribution is $(c_m/2^m)\, N \tau^m$. Every even order
therefore carries the same single power of $N$, with
\begin{equation}
  b_2 = -1, \quad b_4 = -2, \quad b_6 = -\tfrac{437}{48},
  \quad b_8 = -\tfrac{30833}{576}, \quad b_{10} = -\tfrac{3099613}{8640},
  \qquad b_m = c_m/2^m .
\end{equation}
\end{theorem}
\gid{RESULT:SINGLE\_FACE\_PLANAR\_GRADING}

\noindent Recorded \texttt{proven} with \texttt{output-certified}
evidence. The exponent and the coefficients are exact limits of exact
rational functions of $N$, not numerical extrapolations. An independent
integer-rank engine reproduces $c_2, c_4, c_6, c_8$ by extrapolation to
within $1.4 \times 10^{-5}$ relative --- a cross-check by a different
method, not the source of the values.

The word \emph{leading} in Theorem~\ref{thm:grading} is doing work and
should not be dropped. Each $S_m$ is an exact rational function of $N$
with relative corrections in $N^{-2}$, so the reduction to $\tau$ is a
planar-limit statement, not an identity at finite rank. Explicitly, at
second order
\begin{equation}
  S_2(N) = -4N^{-3} - 4N^{-7},
  \qquad\text{so}\qquad
  \frac{S_2\,u^2}{N} = -\tau^2 - \frac{\tau^2}{N^4}.
\end{equation}
The single-face band per unit $N$ is therefore a function of $\tau$
alone \emph{as} $N \to \infty$ at fixed $\tau$, and not at finite $N$.
A statement of the second form would be false, and an earlier internal
summary of this result made it.

\subsection{The grading is a cancellation theorem}

Equation~\eqref{eq:grading} is a much smaller number than its
constituents, and the reason is exact.

\begin{theorem}[Four-channel cancellation]
\label{thm:cancellation}
Decomposing $S_m$ by the intermediate state of the perturbed vector
gives exactly four contributing channels at every even $m$ from $4$ to
$10$: the singlet $\big((),()\big)$, the adjoint
$\big((1),(1)\big)$, and the two-box states $\big((2),()\big)$ and
$\big((1,1),()\big)$. Each is of order $N^{-(m-1)}$, with leading
coefficients
\begin{equation}
  (-4,\ +2,\ +1,\ +1) \times 2^{\,m/2-2},
\end{equation}
which \emph{sum to zero}. At $m = 2$ only the singlet and the adjoint
contribute, at $-2$ and $+2$, and these also cancel.
\end{theorem}
\gid{RESULT:SINGLE\_FACE\_PLANAR\_GRADING}

So the individual channels are of order $N^{-(m-1)}$ and their leading
behaviour annihilates, leaving \eqref{eq:grading} at $N^{-(2m-1)}$ ---
two further orders down in the $N^{-2}$ expansion. At $m = 4$ this is
the statement that channels of order $N^{-3}$ conspire to produce a band
of order $N^{-7}$.

This is why the grading is stated as a theorem about a mechanism rather
than as an observed power. A power law read off from fitted data at
several ranks would carry no information about \emph{why} the leading
terms are absent, and would give no reason to expect
\eqref{eq:grading} to persist. The cancellation does.

\begin{scopenote}
Theorem~\ref{thm:grading} concerns the \emph{one-plaquette} diagonal
splitting at even orders $m \le 10$. It is not the cluster cumulant of
the multi-face conjecture, which is not discharged by it: the one-face
case at sixth order was already established separately, and the
shared-link pair route carries the multi-face statement. The general-$m$
law remains a conjecture. All statements are about coefficients at fixed
order, not about convergence.
\end{scopenote}

\subsection{Odd orders are determinant families}

\begin{theorem}[Centre parity]
\label{thm:centreparity}
For \emph{even} $N$, every odd-order matrix element of the des Cloizeaux
effective operator of the strong-coupling band vanishes identically ---
on every cluster, in both $C$-parity sectors, at every order.

For \emph{odd} $N$, an odd order $m$ vanishes unless $m \ge N-2$. At
$m = N-2$ the only surviving element is the one-plaquette vertex
\begin{equation}
  H_{N-2}(\bar F, F)
  \;=\; -\,\frac{\big(N/(N+1)\big)^{N-3}}{\big((N-3)!\big)^2},
  \label{eq:centrevertex}
\end{equation}
which is $-1$ at $N = 3$ and $-\tfrac{25}{144}$ at $N = 5$. Hops and
leakages vanish at that order.
\end{theorem}
\gid{RESULT:ODD\_ORDER\_CENTRE\_PARITY}

The mechanism is the centre. The $\SU(N)$ Haar integral of a word
vanishes unless every link's net flux is $0 \bmod N$; on a single
plaquette the four links are private, so the intermediate states are
characters with energy $2C_2(R)$, and the flux condition at odd order
can be met only once enough boxes have accumulated to close modulo $N$.
That threshold is $N-2$, and it is why the first odd order is a
function of the rank rather than a fixed number.

\begin{scopenote}
Statements about coefficients of the strong-coupling series at fixed
order, not about convergence. On a periodic lattice the link-set lemma
requires even extents. The even orders and the overlap theorem are not
addressed here.
\end{scopenote}

\subsection{Why this refuted a claim}

The register records a claim that centre-only circuits are dynamically
dark --- that configurations carrying only centre flux cannot contribute
to the dynamics. Theorem~\ref{thm:centreparity} and its fifth-order
companion refute it: the exact fifth-order determinant correction
$\delta c_{5,\mathrm{det}}$ is nonzero. The claim was falsified by an
exact computation rather than by an estimate, which is the only way a
claim of that shape can be settled.
